Generative recommenders depend heavily on the quality of the item tokenizer that maps each item to a discrete semantic ID. Recent collaborative tokenizer designs usually inject one globally fused collaborative representation into the entire tokenizer, but our study of MiniOneRec suggests that this framing is too coarse for hierarchical semantic IDs. We find that the dominant tokenizer bottleneck is not full-code collision after generation, but local ambiguity inside semantically plausible prefixes, where the prefix is correct while the final leaf token remains hard to predict. Based on this observation, we propose \textbf{MGR-SID}, a hierarchy-aware graph-regularized tokenizer built on top of MiniOneRec's semantic \rqvae{} backbone. The current \textbf{v1} instantiation uses a train-only graph bank with a purified coarse collaborative graph, a spectral mid-resolution graph, and a purified local transition graph, and applies different graph regularizers to different SID levels. We first show on Industrial and Office diagnostic probes that spectral middle-resolution graph views are much stronger than heuristic middle views for same-prefix disambiguation. We then align the tokenizer training regime with the upstream MiniOneRec recipe and evaluate training-time graph regularization under this fair setting. On Industrial, the hierarchy-aware variant reduces the final generated index collision from $0.00353$ to $0.00326$ relative to the semantic baseline, and decreases the fraction of test targets in multi-leaf same-$l_2$ buckets from $68.94\%$ to $61.31\%$. Pushing the improved SID into downstream SFT yields a more nuanced result: the hierarchy-aware tokenizer improves short-range ranking ($\mathrm{NDCG}@3: 0.0777 \rightarrow 0.0802$) and lowers same-prefix evaluate errors, but it does not yet surpass the semantic baseline on broader $@10+$ metrics. These results suggest that graph-structured collaboration is most useful when it enters SID learning as level-specific structural supervision rather than uniform feature fusion, while the tokenizer-to-SFT transfer remains only partially solved.
